% ===== §12 The perception of relational value =====
\section{The Perception of Relational Value}
\label{p09:sec:perception}

\noindent
The eudaimonics of the preceding section named a field and stated a first principle, and then honestly marked how much it left undone. It opened a cluster of practical questions---how the field is cultivated, how value is perceived, how value is co-created, how return is secured---and declared them the work of a sequel. This section begins that sequel, and it begins where it must: with \emph{perception}. For one cannot cultivate a field for the circulation of real value, nor secure the return of value to its creators, nor even know whether anyone has been reduced to fuel, without first being able to \emph{perceive} what value is being generated. The whole of the practical work rests on a prior act of discernment, and that act is harder, and stranger, than it first appears.

The difficulty is not incidental but structural, and it follows directly from the theory of value of \xref{p09:sec:value}. Of the three strata, two announce themselves loudly and one is, by its nature, almost inaudible. \emb{Imaginary value and symbolic value are conspicuous; real value is, in the ordinary case, imperceptible.} The labour of perceiving relational value is therefore not a passive registering of something given but an active, disciplined attention that must work \emph{against} the grain of what most readily presents itself.

\subsection{Why the two loud values drown out the third}
\label{p09:sec:loud-values}

Consider why the imaginary and the symbolic are loud. \emb{Imaginary value is loud because it is intense.} The idealization of the other, the felt completeness of the merged dyad, the resplendence of the rose---these flood perception with affect; they are, phenomenologically, almost all one can see when they are present. The lover in the grip of imaginary value does not perceive the real other at all; he perceives the radiant image with which he has filled the void, and the intensity of that image is precisely what blinds him to the irreducible person behind it. Imaginary value does not merely coexist with the perception of real value; it actively occludes it, for the image fills exactly the void that real value would require to be kept open.

\emb{Symbolic value is loud because it is countable.} It comes pre-measured: the price of the gift, the rank of the match, the number of years, the visible markers of a successful relation. The countable is cognitively cheap to perceive---it asks nothing of the perceiver but to read off a magnitude---and it is socially reinforced, for the symbolic is the stratum the surrounding world can see and judge. So perception, left to its own economy, drifts toward the symbolic as water finds the level: it attends to what can be counted because counting is easy and counted things are what others recognize.

Between the intense and the countable, the third stratum has no native loudness at all. \emb{Real value is quiet because it is, by definition, that which neither idealization can image nor measure can count.} It is the shared silence, the bodily presence, the daily mutual witness, the accompaniment of an unfillable void---and none of these announces itself. They are not intense (the calm presence is the opposite of the idealizing flood) and they are not countable (the witness of a wound commensurates into no magnitude). Real value is generated, as \xref{p09:sec:value} argued, only in the encircling interpretive cycle; and a thing generated only in a slow encircling is not the kind of thing that presents itself at a glance. It must be perceived along a path, which is to say it must be \emph{walked} to be seen---the perceptual correlate of the claim that depth is the area a path encloses.

\subsection{Perception as a labour against the grain}
\label{p09:sec:perception-labour}

From this the central claim of the section follows.

\begin{claim}[The perception of real value is a disciplined attention against the grain]
To perceive the real value being generated in a relation is not to register something given but to perform a disciplined attention that works against the grain of perception's own economy---against the intensity of the imaginary, which floods the field and fills the void; and against the easy countability of the symbolic, toward which attention drifts of itself. It is the deliberate keeping-open of the very void that the imaginary rushes to fill, and the deliberate refusal of the measure toward which the symbolic invites one to lapse. Real value is seen only by an attention that declines both the radiant image and the legible number, and consents to the slow, uncounted encircling along which alone the quiet third value comes into view.
\end{claim}

This recasts, at the level of perception, the epistemology of \xref{p09:sec:epistemology}. There we found that the sign of the curvature could be read only in a sustained, fallible, first-personal attention, and never off an external metric. We can now say why, from the side of value: because the value that matters most---the real value that alone bears the good's phase from within---is precisely the value no metric reaches and no image displays. The unobservability of the curvature and the inaudibility of real value are the same fact seen twice. \emb{The labour of diagnostic attention just is the labour of perceiving real value}; to attend to a relation rightly is to keep declining the loud two values long enough for the quiet third to be heard.

And this labour is, in its structure, the same witness the poem demands (\xref{p09:sec:poem}). To perceive real value is not to decode a meaning but to accompany a person along the path on which their irreducible value is generated---to witness, not to read off. The perception of relational value is therefore not a faculty one possesses but a practice one sustains; it is of a piece with the joint attention of the series' earlier papers, and with the non-possessive disposition that lets the other's void stay open rather than filling it with an image of one's own.

\subsection{The two characteristic failures of perception}
\label{p09:sec:perception-failures}

If perception can be done well, it can be done badly, and its failures are not random but track the two loud values. Naming them gives the practice its negative discipline---what it must refuse.

\emc{The imaginary failure of perception} is to mistake the radiant image for the person: to perceive, with great intensity, a value that is in truth one's own projection, and to be unable to perceive the real other at all because the void where that other's irreducibility lives has been filled with the image. This is the perceptual face of ``encircling the idol'' (\xref{p09:sec:non-possess}); and it is, characteristically, the failure of passionate love at its height---the rose perceived in its resplendence, the real person it overgrows unseen. Its correction is not to feel less but to keep the void open: to let the intensity subside enough that the real other, who was never the image, can come into view.

\emc{The symbolic failure of perception} is to mistake the countable for the whole: to perceive only what can be measured---the provision, the status, the visible markers---and to be blind to the uncounted real value being generated or destroyed beneath them. This is the perceptual face of the alienation of social reproduction (\xref{p09:sec:coupled}): the care-labour that creates the value sustaining a relation goes \emph{unperceived} precisely because it is uncountable, and what is unperceived cannot be returned. \emb{The non-perception of real value is the first moment of its non-return.} One cannot return to a creator a value one never saw her create; the reduction of a participant to fuel begins in the failure to perceive the value she generates. The correction is to learn to perceive the uncounted---to attend to the silence, the presence, the witness, that no magnitude records.

This last point binds perception to justice and shows why this section had to come first in the practical sequence. The return of value to its creators (\xref{p09:sec:coupled}) is impossible without the perception of the value created; the diagnosis of whether anyone has been reduced to fuel depends on perceiving the uncounted value that the fuel-position generates and does not receive back. \emb{Perception is the first practical act, because every later one---cultivation, co-creation, return---presupposes that one can see what is being made.} A field cannot be cultivated for a value one cannot perceive; a circulation cannot be made just toward a creation one cannot see. The practice of sustainability begins, then, not in action but in a discipline of attention: the slow, against-the-grain, non-possessive perceiving of the quiet value that neither floods nor counts---and which, unperceived, can neither be cultivated nor returned.
